\section{Adjunction of a Unit Element}

Let A be a k-algebra. In III, §1, No.~2, p.~431, we defined the unital algebra \(\widetilde{A}\) deduced from A by the adjunction of a unit element e. We identify A with a two-sided ideal of \(\widetilde{A}\) . The k-module \(\widetilde{A}\) is the direct sum of the submodules ke and A.

PROPOSITION 3. --- a) Let \(\widetilde{\mathfrak{a}}\) be a left ideal of \(\widetilde{\mathrm{A}}\) such that \(\widetilde{\mathrm{A}} = \widetilde{\mathfrak{a}} + \mathrm{A}\) . We set \(\mathfrak{a} = \widetilde{\mathfrak{a}} \cap \mathrm{A}\) . There exists an element \(u\) of \(\mathrm{A}\) such that \(u - e\) belongs to \(\widetilde{\mathfrak{a}}\) . If \(u\) is such an element, then it is a right unit of \(\mathrm{A}\) modulo \(\mathfrak{a}\) , and we have \(\widetilde{\mathfrak{a}} = \mathfrak{a} + k(u - e)\) ; in particular, the ideal \(\mathfrak{a}\) is regular.

\begin{enumerate}
\def\labelenumi{\alph{enumi})}
\setcounter{enumi}{1}
\item
  Conversely, let \(\mathfrak{a}\) be a regular left ideal of A and \(u\) a right unit of A modulo \(\mathfrak{a}\) . Set \(\widetilde{\mathfrak{a}} = \mathfrak{a} + k(u - e)\) . Then \(\widetilde{\mathfrak{a}}\) is a left ideal of \(\widetilde{\mathrm{A}}\) such that \(\widetilde{\mathrm{A}} = \widetilde{\mathfrak{a}} + \mathrm{A}\) , and we have \(\mathfrak{a} = \widetilde{\mathfrak{a}} \cap \mathrm{A}\) .
\item
  If \(k\) is a field, then the condition \(\widetilde{\mathrm{A}} = \widetilde{\mathfrak{a}} + \mathrm{A}\) is equivalent to saying that \(\widetilde{\mathfrak{a}}\) is not contained in \(\mathrm{A}\) .
\end{enumerate}

Under the assumption of a), the element \(e\) of \(\widetilde{\mathbf{A}}\) can be written as \((e - u) + u\) with \(u \in \mathbf{A}\) and \(u - e \in \widetilde{\mathfrak{a}}\) . Let \(x = \lambda e + a\) be an element of \(\widetilde{\mathbf{A}}\) , where \(\lambda \in k\) and \(a \in \mathbf{A}\) ; if \(x\) belongs to \(\widetilde{\mathfrak{a}}\) , then the element \(y = a + \lambda u = x + \lambda (u - e)\) of \(\mathbf{A}\) belongs to \(\mathfrak{a}\) , and we have \(x = y - \lambda (u - e)\) ; this proves the equality \(\widetilde{\mathfrak{a}} = \mathfrak{a} + k(u - e)\) . Moreover, for every \(a\) in \(\mathbf{A}\) , the element \(au - a = a(u - e)\) of \(\mathbf{A}\) belongs to \(\mathfrak{a}\) , so \(u\) is a right unit modulo \(\mathfrak{a}\) . This proves assertion a).

Let \(\mathfrak{a}\) and \(u\) be as in b). Since \(\widetilde{\mathrm{A}}\) is the direct sum of \(\mathrm{A}\) and \(k(u - e)\) , we have \(\widetilde{\mathfrak{a}} + \mathrm{A} = \widetilde{\mathrm{A}}\) and \(\widetilde{\mathfrak{a}} \cap \mathrm{A} = \mathfrak{a}\) . Every element of \(\widetilde{\mathfrak{a}}\) is of the form \(x + \lambda (u - e)\) with \(x \in \mathfrak{a}\) and \(\lambda \in k\) . For every \(a \in \mathrm{A}\) , we have

\[
a (x + \lambda (u - e)) = a x + \lambda (a u - a).
\]

Since u is a right unit of A modulo \(\mathfrak{a}\), the sum \(ax + \lambda(au - a)\) belongs to \(\mathfrak{a}\), so that \(\widetilde{\mathfrak{a}}\) is a left ideal of A. This proves b).

Finally, if \(k\) is a field, then \(A\) is a hyperplane in \(\widetilde{A}\) , and \(\widetilde{A}\) is the sum of \(\widetilde{\mathfrak{a}}\) and \(A\) if and only if \(\widetilde{\mathfrak{a}}\) is not contained in \(A\) .

COROLLARY. --- The regular left ideals of A are the ideals of the form \(A \cap \widetilde{\mathfrak{a}}\) , where \(\widetilde{\mathfrak{a}}\) is a left ideal of \(\widetilde{A}\) such that \(\widetilde{A} = \widetilde{\mathfrak{a}} + A\) .

PROPOSITION 4. --- a) Let \(\mathfrak{a}\) be a regular maximal left ideal of A. There exists a unique left ideal \(\widetilde{\mathfrak{a}}\) of \(\widetilde{\mathrm{A}}\) such that \(\widetilde{\mathrm{A}} = \widetilde{\mathfrak{a}} + \mathrm{A}\) and \(\mathfrak{a} = \widetilde{\mathfrak{a}} \cap \mathrm{A}\) . This ideal is maximal and does not contain A.

\begin{enumerate}
\def\labelenumi{\alph{enumi})}
\setcounter{enumi}{1}
\tightlist
\item
  The mapping \(\widetilde{\mathfrak{a}}\mapsto \widetilde{\mathfrak{a}}\cap \mathrm{A}\) is a bijection from the set of maximal left ideals of \(\widetilde{\mathrm{A}}\) not containing A to the set of regular maximal left ideals of A.
\end{enumerate}

Let \(\mathfrak{a}\) be a regular maximal left ideal of A. By Proposition 3, a), the left ideals \(\mathfrak{b}\) of \(\widetilde{\mathrm{A}}\) such that \(\mathfrak{b} + \mathrm{A} = \widetilde{\mathrm{A}}\) and \(\mathfrak{b} \cap \mathrm{A} = \mathfrak{a}\) are the ideals \(\mathfrak{a} + k(u - e)\) , where \(u\) is a right unit modulo \(\mathfrak{a}\) . To prove the uniqueness of \(\widetilde{\mathfrak{a}}\) , it therefore suffices to prove that two right units \(u\) and \(u'\) of A modulo \(\mathfrak{a}\) are congruent modulo \(\mathfrak{a}\) . Let us reason by contradiction and suppose that \(u - u'\) does not belong to \(\mathfrak{a}\) . The formula \(x(u - u') = (xu - x) - (xu' - x)\) shows that we have \(\mathrm{A}(u - u') \subset \mathfrak{a}\) ; it follows that \(\mathfrak{a} + k(u - u')\) is a left ideal of A containing \(\mathfrak{a}\) and distinct from \(\mathfrak{a}\) . Since \(\mathfrak{a}\) is maximal, we therefore have \(\mathfrak{a} + k(u - u') = \mathrm{A}\) , and so \(\mathrm{AA} \subset \mathfrak{a}\) . For every \(x \in \mathrm{A}\) , we have \(x \equiv xu (\mathrm{mod} \mathfrak{a})\) , and therefore \(x \in \mathfrak{a}\) by the above, which contradicts the assumption \(\mathfrak{a} \neq \mathrm{A}\) .

Hence there exists a unique left ideal \(\widetilde{\mathfrak{a}}\) of \(\widetilde{\mathrm{A}}\) such that \(\widetilde{\mathrm{A}} = \widetilde{\mathfrak{a}} + \mathrm{A}\) and \(\mathfrak{a} = \widetilde{\mathfrak{a}} \cap \mathrm{A}\) . Let \(\mathfrak{b}\) be a proper left ideal of \(\widetilde{\mathrm{A}}\) containing \(\widetilde{\mathfrak{a}}\) . Then \(\mathfrak{b} \cap \mathrm{A}\) is a proper left ideal of \(\mathrm{A}\) containing \(\mathfrak{a}\) . It is therefore equal to \(\mathfrak{a}\) because \(\mathfrak{a}\) is maximal, which implies \(\mathfrak{b} = \widetilde{\mathfrak{a}}\) by Lemma 2 of VIII, p.~4. This proves that \(\widetilde{\mathfrak{a}}\) is a maximal ideal of \(\widetilde{\mathrm{A}}\) ; for such an ideal, the condition \(\widetilde{\mathrm{A}} = \widetilde{\mathfrak{a}} + \mathrm{A}\) means that \(\widetilde{\mathfrak{a}}\) does not contain A.

It remains to prove that if \(\widetilde{\mathfrak{a}}\) is a maximal left ideal of \(\tilde{A}\) that does not contain A, then the left ideal \(\mathfrak{a} = A \cap \widetilde{\mathfrak{a}}\) of A is maximal. Let \(\mathfrak{b}\) be a proper left ideal of A containing \(\mathfrak{a}\). Let u be a right unit modulo \(\mathfrak{a}\) such that \(\widetilde{\mathfrak{a}} = \mathfrak{a} + k(u - e)\) (Proposition 3). It is also a right unit modulo \(\mathfrak{b}\); denote by \(\widetilde{\mathfrak{b}}\) the ideal \(\mathfrak{b} + k(u - e)\) of \(\tilde{A}\) . By Proposition 3, b), we have \(\mathfrak{b} = A \cap \widetilde{\mathfrak{b}}\) . The ideal \(\widetilde{\mathfrak{a}}\) , equal to \(\mathfrak{a} + k(u - e)\) , is contained in \(\widetilde{\mathfrak{b}}\) and is therefore equal to \(\widetilde{\mathfrak{b}}\) because \(\widetilde{\mathfrak{b}}\) is distinct from \(\tilde{A}\) and \(\widetilde{\mathfrak{a}}\) is maximal. Consequently, we have \(\mathfrak{a} = \mathfrak{b}\) , and \(\mathfrak{a}\) is maximal.

We leave it to the reader to translate Propositions 3 and 4 to right ideals.

